\chapter{Trabalhos relacionados}
\label{cap:relacionados}

% Fontes: wiki/sources/*, comparisons/Tool Capability Matrix,
% comparisons/LLM Integration Patterns, benchmarks/*.
% Todo número é o que os autores reportam, não uma medição nossa.

Este capítulo apresenta duas literaturas. A primeira
(Seção~\ref{sec:rel-interrupcoes}) trata da detecção de defeitos de
concorrência em programas orientados a interrupções. A segunda
(Seção~\ref{sec:rel-llm}) trata do uso de modelos de linguagem para apoiar a
análise estática. A Seção~\ref{sec:rel-benchmarks} descreve os benchmarks da primeira literatura, e a Seção~\ref{sec:rel-comparacao} compara os trabalhos e
formaliza a lacuna que este trabalho ocupa.

Salvo indicação em contrário, os números citados são os relatados pelos
próprios autores. Muitos deles foram medidos em hardware diferente, sobre
contagens diferentes do mesmo benchmark, ou copiados de outros artigos porque
as ferramentas de comparação não estavam disponíveis; por isso não devem ser
comparados diretamente entre si.

\section{Detecção de defeitos em programas com interrupções}
\label{sec:rel-interrupcoes}

\subsection{SDRacer}

\citeonline{sdracer} combinam análise estática, execução simbólica e simulação
numa plataforma virtual para detectar, validar e reparar condições de corrida
em software embarcado. O trabalho é o que melhor enuncia por que as técnicas de
\textit{threads} não se transferem: interrupções não bloqueiam, a preempção é
assimétrica, a sincronização é por desabilitação, a ocorrência depende do
estado do hardware e o reparo exige primitivas específicas de interrupção.

A ferramenta tem quatro etapas. A análise estática emite alertas de corrida; a
execução simbólica gera entradas e intercalações que alcançam cada alerta e
descarta os de condições de caminho conflitantes; a plataforma virtual Simics
força as interrupções a ocorrer nos pontos suspeitos e mantém apenas os pares
que de fato se intercalam; e, por fim, a ferramenta sugere correções: desabilitar e reabilitar a interrupção em volta do acesso, inserir travas, estender uma seção crítica existente ou fundir seções adjacentes.

Os autores relatam 190 condições de corrida detectadas em nove benchmarks
embarcados. A execução simbólica eliminou 40,3\% dos alertas estáticos e a
validação dinâmica, mais 36,7\%. As correções custaram menos de 0,09 de
sobrecarga em 9 dos 11 programas avaliados; em dois, a sobrecarga foi bem
maior, porque desabilitar interrupções alterou o fluxo de controle da tarefa
principal. A ausência de falsos negativos foi afirmada por inspeção manual. É
a única ferramenta revisada que repara defeitos, e sua detecção se limita ao
que a plataforma virtual consegue simular.

\subsection{IntRace}

\citeonline{intrace} detectam \textit{data races} de forma puramente estática,
em três etapas de custo e precisão crescentes. A primeira constrói o CFG a
partir do compilador Clang e de uma passagem LLVM, identifica os recursos
compartilhados por análise de ponteiros interprocedural e casa pares de
acessos, ignorando deliberadamente condições de desvio e estado das
interrupções, para não perder nada. A segunda descarta pares cujos fluxos não
podem se sobrepor, por prioridade ou mascaramento; mecanismos de mascaramento
não padronizados precisam ser informados pelo usuário num arquivo de
configuração. A terceira verifica a viabilidade de cada par restante com o
solver Z3.

A medição por etapa é o resultado mais aproveitável do artigo: cerca de 208
candidatos por programa real após a primeira etapa; a segunda elimina 54,2\%;
a terceira elimina mais 86,2\% do que lhe chega e consome 69,7\% do tempo total
de análise. Sobre 30 dos 31 casos simples do Racebench, os autores relatam 5 falsos positivos e 73,2\% a
menos de falsos positivos que a ferramenta Rchecker, cujos números, porém,
foram copiados do artigo original, pois a ferramenta não estava disponível.
Sobre nove programas industriais, relatam 118 condições de corrida e 9 falsos
positivos. As duas causas de falsos positivos que os autores reconhecem são
operações implícitas de interrupção e condições de corrida benignas deixadas de
propósito.

\subsection{NIChecker}

\citeonline{nichecker} verificam violações de atomicidade traduzindo o programa
com interrupções num programa sequencial, pela técnica de sequencialização
preguiçosa, e entregando-o ao verificador CBMC \cite{lazycseq,cbmc}. As
\textit{condições de invocação} de cada \isr\ no programa sequencial codificam
prioridade e mascaramento, o que dá suporte nativo ao aninhamento. Antes do
desenrolamento, a ferramenta aplica abstração de laços, o que lhe permite
encontrar defeitos que exigiriam desenrolar um laço 10.000 vezes; o fatiamento
de programa e a redução de pontos de preempção aceleram a verificação em até
42,2\%. É o único trabalho revisado que prova a correção limitada de sua
tradução.

Os autores relatam 96,4\% de precisão no Racebench e 47 violações encontradas,
sem falsos positivos, numa suíte de 18 programas reais. A principal limitação
prática é que o usuário precisa indicar, a cada execução, qual variável global
e qual padrão verificar.

\subsection{BMC4AV}

\citeonline{bmc4av} também usam verificação limitada sobre o CBMC, mas em vez
de sequencializar o programa guiam o solver com um grafo de acessos à memória
construído durante a resolução. Primeiro identificam violações
\textit{potenciais} por casamento de padrões e, para cada uma, as relações de
ordem que precisariam valer; depois estendem o grafo guiados por essas relações
até \textit{confirmar} ou não a violação. A ablação dos próprios autores, feita
numa só máquina, é o resultado mais confiável do trabalho: a orientação pelas
relações de ordem dá a precisão, e o grafo dá a eficiência.

Os autores afirmam encontrar 94 de 94 violações na suíte de 18 programas reais
e afirmam que o NIChecker encontra apenas 37 delas, com 57 falsos negativos e taxa de acerto de 39,4\%.
Trata-se de um \textit{preprint} sem revisão por pares, e a contagem de
94 usa uma unidade diferente da usada pelo NIChecker
(Seção~\ref{sec:rel-comparacao}). Os autores ainda classificam o padrão
$(R,W,W)$ como benigno e excluem do Racebench os seis casos que só o contêm.

\subsection{Outras ferramentas citadas}

Os trabalhos acima comparam-se com outras ferramentas: o Rchecker, detector de
condições de corrida baseado no CBMC com uma lista de máscaras de interrupção;
o intAtom, detector estático de violações de atomicidade; o CPA4AV, baseado em
árvores de alcançabilidade abstratas, que falha nos casos com laços profundos;
e o iCBMC, extensão do CBMC para interrupções aninhadas. O Rchecker e o intAtom não puderam ser obtidos pelos autores que os usaram como comparação, e seus números foram copiados de outros artigos, medidos em outro hardware \cite{intrace,nichecker,bmc4av}.

\section{Modelos de linguagem aplicados à análise estática}
\label{sec:rel-llm}

\subsection{Triagem de alertas}

\citeonline{llift} partem do UBITect, analisador que procura usos de variáveis
antes da inicialização no \textit{kernel} Linux. A análise estática do UBITect
gera cerca de 140 mil candidatos, e sua etapa de execução simbólica abandona
53 mil deles (40\%) por limite de tempo ou memória. O LLift entrega apenas esses
casos não decididos a um \llm, com uma política conservadora: tudo o que não
for provado seguro é reportado. O trabalho encontrou quatro defeitos reais
confirmados pela comunidade do Linux. Sua ablação mostra o recall passando de
0,15, com uma pergunta direta, para 1,00, com os quatro componentes de
\textit{prompt} descritos na Seção~\ref{sec:llm}. É a referência estrutural
mais próxima deste trabalho, e seu formato de ablação é adotado aqui.

\citeonline{chatgpt-static} avaliam modelos GPT na remoção de falsos
positivos do analisador Infer e obtêm precisão de 93,88\% para desreferência de
ponteiro nulo, mas apenas 63,33\% para vazamento de recursos, com os mesmos
\textit{prompts} e o mesmo modelo. A lição é que a qualidade da triagem depende
da classe de defeito e não se transfere por suposição.

\citeonline{reducing-fp} filtram alertas do analisador CppCheck com um modelo
CodeBERT ajustado para prever funções vulneráveis. No conjunto Big-Vul, a
precisão sobe de 5,88\% para 96,20\%, mas o recall cai de 46,90\% para
40,66\%: os autores reconhecem que a filtragem introduz falsos negativos. É
desse trabalho que vem a métrica \textit{Inspection Ratio}, que cai de
50,49\% para 2,44\%.

\subsection{Inferência de especificações}

\citeonline{iris} usam um \llm\ para inferir quais funções de bibliotecas são
fontes e sumidouros de dados contaminados, alimentando com essas
especificações o analisador CodeQL. Num conjunto de 120 vulnerabilidades reais
em Java, a detecção sobe de 27, com o CodeQL sozinho, para 55. A precisão
continua baixa (taxa média de falsas descobertas de 84,82\%), e a triagem por
\llm\ só melhora a precisão com modelos grandes; com modelos pequenos, piora.

\citeonline{adataint} seguem a mesma ideia, combinando inferência de
especificações com um classificador que filtra alertas. Os resultados vêm de
benchmarks sintéticos (Juliet e SV-COMP), e o filtro descarta alertas, com
recall de 75,4\%, isto é, cerca de um quarto dos defeitos é perdido por
construção.

\subsection{Detecção e reparo}

\citeonline{skipanalyzer} encadeiam três componentes com \llm: detecção,
filtragem de falsos positivos e reparo, sobre alertas do Infer em Java. No
reparo, obtêm taxa lógica (correção equivalente à escrita por um humano) de
97,30\% para desreferência nula e 91,77\% para vazamento de recursos, sem
ajuste fino. O mesmo grupo é autor de \citeonline{chatgpt-static}, e os dois
trabalhos compartilham dados; devem ser lidos como um só conjunto de evidências.

\citeonline{sastgenius} integram o analisador Semgrep a um modelo Llama~3
ajustado, que faz triagem, geração de \textit{exploits} para validar os
alertas e sugestão de correções. Trata-se de um \textit{preprint}, que
relata a precisão subindo de 35,7\% para 89,5\% mas não relata recall. Sua
contribuição mais relevante aqui é apontar o risco de \textit{prompt
injection}: o código analisado, incluindo comentários, entra no
\textit{prompt} e pode tentar instruir o modelo a ignorar um defeito.

\section{Benchmarks}
\label{sec:rel-benchmarks}

\subsection{Racebench}

O Racebench \cite{racebench} é o benchmark sobre o qual o IntRace, o NIChecker e o BMC4AV são avaliados.
Foi criado para uma competição de protótipos de detecção de condições de
corrida realizada em Hangzhou, na China, em novembro de 2019, e seus programas
são modelados em software embarcado do setor aeroespacial. A versão 2.1 tem 31
casos simples e 2 complexos. Cada programa tem um fluxo principal e de uma a
três \isr s, com até três níveis de prioridade.

O gabarito é anotado no próprio código, e cada defeito é uma
tripla de acessos, o que significa que o benchmark especifica
violações de atomicidade mesmo para os trabalhos que o usam para
\textit{data races}. Além dos defeitos, o benchmark contém armadilhas
(\textit{traps}): falsos positivos plantados de propósito, vários deles trechos corretamente protegidos por mascaramento. Na competição, nos casos simples, cada detecção valia 2 pontos e cada falso positivo custava 3, uma regra que premia a
precisão muito acima do recall.

Contando diretamente as anotações dos 31 casos simples, este trabalho obteve
48 defeitos e 38 armadilhas. Os artigos usam outras contagens para os
mesmos casos: 50 \cite{intrace}, 54 \cite{nichecker} e 38 sobre um subconjunto de
25 casos \cite{bmc4av}. As anotações usam ao menos quatro gramáticas
incompatíveis e contêm erros de digitação: um leitor estrito encontra só 28
defeitos, sem acusar erro. Essa é uma causa plausível de parte das divergências
entre os artigos.

\subsection{Suíte de programas reais}

A segunda suíte, usada pelo NIChecker e pelo BMC4AV, tem 18 programas extraídos
de seis pacotes: \textit{firmware} de um registrador de temperatura, controle de
LEDs, um sistema de freio eletrônico gerado a partir de um modelo Simulink e
três \textit{drivers} (barramento I$^2$C e dois temporizadores de
\textit{watchdog}, ao menos um deles extraído do \textit{kernel} Linux), num
total de 10.633 linhas \cite{nichecker,bmc4av}. Cada programa acompanha um arquivo com as
prioridades dos fluxos e um arquivo de gabarito, que contém 45 entradas de
violação verdadeira e 14 de falso positivo.

\section{Comparação e lacuna}
\label{sec:rel-comparacao}

A Tabela~\ref{tab:capacidades} compara as quatro ferramentas de detecção.

\begin{table}[htb]
  \centering
  \caption{Capacidades das ferramentas de detecção revisadas}
  \label{tab:capacidades}
  \small
  \begin{tabular}{@{}L{2.8cm}L{2.6cm}L{2.6cm}L{3.2cm}L{2.7cm}@{}}
    \toprule
    & SDRacer & IntRace & NIChecker & BMC4AV \\
    \midrule
    Ano, veículo & 2020, TSE & 2024, J. Supercomp. & 2025, TOSEM & 2026, \textit{preprint} \\
    Classe de defeito & \textit{data race} & \textit{data race} & violação de atomicidade & violação de atomicidade \\
    Método & estático, simbólico e dinâmico & estático e Z3 & sequencialização e BMC & BMC com grafo de acessos \\
    Aninhamento & exclui reentrância & conversão sequencial & nativo & por ordens parciais \\
    Automação & automática & arquivo de configuração & variável e padrão por execução & automática, segundo os autores \\
    Repara defeitos & sim & não & não & não \\
    Código disponível & não & não & não & sim \\
    \bottomrule
  \end{tabular}
  \fonte{Elaborada pelos autores a partir de \citeonline{sdracer},
  \citeonline{intrace}, \citeonline{nichecker} e \citeonline{bmc4av}.}
\end{table}

Três observações seguem da tabela e da revisão.

Os números não se comparam entre artigos. A disputa mais aguda é a do
BMC4AV contra o NIChecker na suíte real: o NIChecker relata 47 violações, o
BMC4AV credita a ele 37 e afirma que existem 94. O gabarito
distribuído com a suíte registra uma entrada por \textit{variável} compartilhada
(45 entradas), enquanto a tabela do BMC4AV conta uma entrada por
\textit{instância de tripla de acessos}. As duas contagens coincidem nos
programas pequenos e divergem nos grandes, crescendo com o tamanho do programa, o que é explicado mais simplesmente por uma diferença de unidade do que
por 57 defeitos perdidos. Isso não invalida a crítica do BMC4AV, mas mostra que
qualquer taxa de acerto depende de uma escolha de contagem que os artigos não
explicitam.

A ausência de falsos negativos é afirmada, não medida. O SDRacer e o IntRace a apoiam em inspeção manual; os trabalhos de verificação limitada, na ausência de
contraexemplo até um limite. Nenhum declara as premissas sob as quais nenhum
defeito é perdido.

As duas literaturas não se encontram. Nenhuma ferramenta de detecção
trata as duas classes de defeito numa só análise, e só a mais antiga repara. Do
outro lado, os sete trabalhos com \llm\ tratam defeitos sequenciais, e a literatura converge para a mesma divisão de trabalho: o analisador faz o raciocínio sobre o
programa inteiro e o modelo faz o julgamento local. Os resultados também
mostram que a política de decisão importa mais que o modelo: filtros que
descartam perdem de 6 a 25 pontos de recall \cite{reducing-fp,adataint},
enquanto a política conservadora do LLift atinge recall de 1,00 no mesmo ponto
de integração \cite{llift}. Nenhum trabalho aplica essa divisão de trabalho a
defeitos de concorrência com interrupções. Essa é a lacuna ocupada por este
trabalho.
